\subsection{Scope}
\noindent To ensure the feasibility of the aBridgeAI, the scope is constrained across five dimensions:

\begin{itemize}
    \item \textbf{Technological Scope:} The system is developed as a responsive web application to ensure cross-platform compatibility. The stack encompasses React for the frontend, FastAPI (Python) for backend and AI control-plane orchestration, PostgreSQL with pgvector for transactional and vector data, Neo4j for the course knowledge graph, Redis for caching and background jobs, Garage for S3-compatible object storage, and LiveKit for WebRTC interview rooms. The target deployment runs a separate native LiveKit worker for conversational interviews; Google OAuth 2.0 / OpenID Connect provides delegated authentication with in-platform second-factor (TOTP) enrolment and session management.
    
    \item \textbf{Temporal Scope:} The development, testing, and deployment lifecycle of this Minimum Viable Product (MVP) is strictly bounded within Semester 252, commencing in early February and concluding in June.
    
    \item \textbf{Spatial and Contextual Scope:} The system's operational context is designed specifically for academic and structured learning environments. It targets higher education institutions, university departments, and coding bootcamps, focusing exclusively on bridging the gap between theoretical academic curricula and practical industry interview requirements.
    
    \item \textbf{Functional Scope:} The system's core functionalities are limited to four main interactive pillars: 
    \begin{enumerate}
        \item \textit{Smart Content Ingestion:} Parsing uploaded documents (PDF, DOCX) and media (MP4) to construct a localized Course Knowledge Graph.
        \item \textit{The Learning Loop:} Providing automated, adaptive quizzes based on the ingested content.
        \item \textit{The Application Loop:} Executing context-aware mock interviews supporting text-based and asynchronous/WebRTC voice-based interactions.
        \item \textit{Assessment Analytics:} Generating automated Gap Reports and tracking student performance matrices.
    \end{enumerate}
    The following capabilities are explicit non-goals of the MVP and are excluded from this phase: synchronous human-to-human tutoring and live video conferencing; in-platform email/password registration, password reset, ``remember me'' device-trust, and per-device session pinning. Authentication is delegated exclusively to Google OAuth, and self-service password recovery is therefore handled by Google's identity infrastructure rather than in-platform.
    
    \item \textbf{Data Scope:} The system is scoped to process, store, and manage standard educational materials, user profiles, and assessment logs. For AI operations, the system limits retrieval data to vector embeddings (via pgvector) derived from institution-uploaded materials. Live interview audio is normally processed transiently for Speech-to-Text (STT) and Text-to-Speech (TTS). Optional audio recording is separately deployment-enabled, consent-gated by policy version, and retained through its own object-storage lifecycle; recording failure does not alter assessment outcome.
\end{itemize}